% ******************************************************************************
% ****************************** Custom Margin *********************************

% Add `custommargin' in the document class options to use this section
% Set {innerside margin / outerside margin / topmargin / bottom margin}  and
% other page dimensions
\ifsetCustomMargin
  \RequirePackage[left=37mm,right=30mm,top=35mm,bottom=30mm]{geometry}
  \setFancyHdr % To apply fancy header after geometry package is loaded
\fi

% Add spaces between paragraphs
%\setlength{\parskip}{0.5em}
% Ragged bottom avoids extra whitespaces between paragraphs
\raggedbottom
% To remove the excess top spacing for enumeration, list and description
%\usepackage{enumitem}
%\setlist[enumerate,itemize,description]{topsep=0em}

% *****************************************************************************
% ******************* Fonts (like different typewriter fonts etc.)*************

% Add `customfont' in the document class option to use this section

\ifsetCustomFont
  % Set your custom font here and use `customfont' in options. Leave empty to
  % load computer modern font (default LaTeX font).
  %\RequirePackage{helvet}

  % For use with XeLaTeX
  %  \setmainfont[
  %    Path              = ./libertine/opentype/,
  %    Extension         = .otf,
  %    UprightFont = LinLibertine_R,
  %    BoldFont = LinLibertine_RZ, % Linux Libertine O Regular Semibold
  %    ItalicFont = LinLibertine_RI,
  %    BoldItalicFont = LinLibertine_RZI, % Linux Libertine O Regular Semibold Italic
  %  ]
  %  {libertine}
  %  % load font from system font
  %  \newfontfamily\libertinesystemfont{Linux Libertine O}
\fi

% *****************************************************************************
% **************************** Custom Packages ********************************
\usepackage{amsmath,amssymb}
\usepackage{booktabs}
\usepackage{tikz-feynman}
\usepackage{xspace}
\usepackage{graphicx}
\usepackage{threeparttable}
\usepackage{adjustbox}
\usepackage{bm}
% ************************* Algorithms and Pseudocode **************************

%\usepackage{algpseudocode}

% ********************Captions and Hyperreferencing / URL **********************

% Captions: This makes captions of figures use a boldfaced small font.
%\RequirePackage[small,bf]{caption}

\RequirePackage[labelsep=space,tableposition=top]{caption}
\renewcommand{\figurename}{Fig.} %to support older versions of captions.sty


% *************************** Graphics and figures *****************************

%\usepackage{rotating}
%\usepackage{wrapfig}

% Uncomment the following two lines to force Latex to place the figure.
% Use [H] when including graphics. Note 'H' instead of 'h'
%\usepackage{float}
%\restylefloat{figure}

% Subcaption package is also available in the sty folder you can use that by
% uncommenting the following line
% This is for people stuck with older versions of texlive
%\usepackage{sty/caption/subcaption}
\usepackage{subcaption}

% ********************************** Tables ************************************
\usepackage{booktabs} % For professional looking tables
\usepackage{multirow}

%\usepackage{multicol}
%\usepackage{longtable}
%\usepackage{tabularx}


% *********************************** SI Units *********************************
\usepackage{siunitx} % use this package module for SI units


% ******************************* Line Spacing *********************************

% Choose linespacing as appropriate. Default is one-half line spacing as per the
% University guidelines

% \doublespacing
% \onehalfspacing
% \singlespacing


% ************************ Formatting / Footnote *******************************

% Don't break enumeration (etc.) across pages in an ugly manner (default 10000)
%\clubpenalty=500
%\widowpenalty=500

%\usepackage[perpage]{footmisc} %Range of footnote options


% *****************************************************************************
% *************************** Bibliography  and References ********************

%\usepackage{cleveref} %Referencing without need to explicitly state fig /table

% Add `custombib' in the document class option to use this section
% \ifuseCustomBib
   % \RequirePackage[square, sort, numbers, authoryear]{natbib} % CustomBib
% customised reference sorting
\ifuseCustomBib
   \RequirePackage[square, sort&compress, numbers]{natbib} % CustomBib

% If you would like to use biblatex for your reference management, as opposed to the default `natbibpackage` pass the option `custombib` in the document class. Comment out the previous line to make sure you don't load the natbib package. Uncomment the following lines and specify the location of references.bib file

%\RequirePackage[backend=biber, style=numeric-comp, citestyle=numeric, sorting=nty, natbib=true]{biblatex}
%Location of references.bib only for biblatex, Do not omit the .bib extension from the filename.
%\addbibresource{References/ref1}
%\addbibresource{References/references}


\fi

% changes the default name `Bibliography` -> `References'
\renewcommand{\bibname}{References}


% ******************************************************************************
% ************************* User Defined Commands ******************************
% ******************************************************************************


% *********** To change the name of Table of Contents / LOF and LOT ************

%\renewcommand{\contentsname}{My Table of Contents}
%\renewcommand{\listfigurename}{My List of Figures}
%\renewcommand{\listtablename}{My List of Tables}


% ********************** TOC depth and numbering depth *************************

\setcounter{secnumdepth}{2}
\setcounter{tocdepth}{2}


% ******************************* Nomenclature *********************************

% To change the name of the Nomenclature section, uncomment the following line

%\renewcommand{\nomname}{Symbols}


% ********************************* Appendix ***********************************

% The default value of both \appendixtocname and \appendixpagename is `Appendices'. These names can all be changed via:

%\renewcommand{\appendixtocname}{List of appendices}
%\renewcommand{\appendixname}{Appndx}

% *********************** Configure Draft Mode **********************************

% Uncomment to disable figures in `draft'
%\setkeys{Gin}{draft=true}  % set draft to false to enable figures in `draft'

% These options are active only during the draft mode
% Default text is "Draft"
%\SetDraftText{DRAFT}

% Default Watermark location is top. Location (top/bottom)
%\SetDraftWMPosition{bottom}

% Draft Version - default is v1.0
%\SetDraftVersion{v1.1}

% Draft Text grayscale value (should be between 0-black and 1-white)
% Default value is 0.75
%\SetDraftGrayScale{0.8}


% ******************************** Todo Notes **********************************
%% Uncomment the following lines to have todonotes.

%\ifsetDraft
%	\usepackage[colorinlistoftodos]{todonotes}
%	\newcommand{\mynote}[1]{\todo[author=kks32,size=\small,inline,color=green!40]{#1}}
%\else
%	\newcommand{\mynote}[1]{}
%	\newcommand{\listoftodos}{}
%\fi

% Example todo: \mynote{Hey! I have a note}

% *****************************************************************************
% ******************* Better enumeration my MB*************
\usepackage{enumitem}

% ******************* My custom preamble ***************
\newcommand{\subsubtitle}[1]{\noindent\textbf{#1} \quad}
% Common physics notation; commands work in both text and math mode.
\newcommand*{\et}{\ensuremath{E_{\mathrm{T}}}\xspace}
\newcommand*{\met}{\ensuremath{E_{\mathrm{T}}^{\mathrm{miss}}}\xspace}
\newcommand*{\mT}{\ensuremath{m_{\mathrm{T}}}\xspace}
\newcommand*{\pT}{\ensuremath{p_{\mathrm{T}}}\xspace}
\newcommand*{\gU}{\ensuremath{g_U}\xspace}
\newcommand*{\MU}{\ensuremath{M_U}\xspace}
\newcommand*{\MLQ}{\ensuremath{M_{\mathrm{LQ}}}\xspace}
% Common background processes; lepton flavour follows the surrounding text.
\newcommand*{\zll}{\ensuremath{Z\rightarrow\ell\ell}\xspace}
\newcommand*{\zee}{\ensuremath{Z\rightarrow ee}\xspace}
\newcommand*{\zmumu}{\ensuremath{Z\rightarrow\mu\mu}\xspace}
\newcommand*{\ztautau}{\ensuremath{Z\rightarrow\tau\tau}\xspace}
\newcommand*{\znunu}{\ensuremath{Z\rightarrow\nu\nu}\xspace}
\newcommand*{\wlnu}{\ensuremath{W\rightarrow\ell\nu}\xspace}
\newcommand*{\wenu}{\ensuremath{W\rightarrow e\nu}\xspace}
\newcommand*{\wmunu}{\ensuremath{W\rightarrow\mu\nu}\xspace}
\newcommand*{\wtaunu}{\ensuremath{W\rightarrow\tau\nu}\xspace}
\newcommand*{\wjets}{\ensuremath{W+\mathrm{jets}}\xspace}
\newcommand*{\zjets}{\ensuremath{Z+\mathrm{jets}}\xspace}
\newcommand*{\wzjets}{\ensuremath{W/Z+\mathrm{jets}}\xspace}
\newcommand*{\ttbar}{\ensuremath{t\bar{t}}\xspace}
% Compatibility with the original command names.
\newcommand*{\Wtaunu}{\wtaunu}
\newcommand*{\mt}{\mT}
\newcommand*{\pt}{\pT}
\newcommand*{\RD}{\ensuremath{R(D^{(*)})}}
\newcommand{\DphiTauMet}{\ensuremath{\Delta\phi(\tau, \met)}\xspace}
\newcommand{\DphiLepMet}{\ensuremath{\Delta\phi(\ell, \met)}\xspace}
\newcommand{\DphiJetMet}{\ensuremath{\Delta\phi(\met, \pT^{\text{jet}})}\xspace}
\newcommand{\minDphiJetPt}{\ensuremath{\pT^{\text{jet}(\min{\Delta\phi)}}}\xspace}
\newcommand{\betaL}{\ensuremath{\beta_L^{23}}\xspace}
\newcommand{\MeV}{\ensuremath{\mathrm{MeV}}}
\newcommand{\GeV}{\ensuremath{\mathrm{GeV}}}
\newcommand{\TeV}{\ensuremath{\mathrm{TeV}}}
\newcommand{\um}{\ensuremath{\mu\mathrm{m}}}
\newcommand{\mm}{\ensuremath{\mathrm{mm}}}
\newcommand{\fb}{\ensuremath{\mathrm{fb}}}
\newcommand{\ifb}{\ensuremath{\mathrm{fb}^{-1}}}
\usepackage{url}
\DeclareMathAlphabet{\mathcal}{OMS}{cmsy}{m}{n}
\usepackage{CJKutf8}

\newcommand{\SR}{\ensuremath{\mathrm{SR1b}}\xspace}
\newcommand{\bVetoSR}{\ensuremath{\mathrm{SR0b}}\xspace}
\newcommand{\CRW}{\ensuremath{\mathrm{CRW}}\xspace}
\newcommand{\CRTop}{\ensuremath{\mathrm{CRTop}}\xspace}
\newcommand{\VRb}{\ensuremath{\mathrm{VRW}}\xspace}
\newcommand{\VRW}{\ensuremath{\mathrm{SR0b}}\xspace}
\newcommand{\VRTop}{\ensuremath{\mathrm{VRTop}}\xspace}

\newcommand{\Res}{\ensuremath{\mathrm{Res}}\xspace}
\newcommand{\NonRes}{\ensuremath{\mathrm{NonRes}}\xspace}
